% =====================================================================
\section{Metodología de ingeniería de requisitos}
\label{sec:metodologia}
 
\subsection{Marco normativo aplicado}
 
El ERS se estructura conforme a ISO/IEC/IEEE 29148:2018, que define el contenido de una
especificación de requisitos de software y las características exigibles a un requisito individual y
al conjunto \cite{iso29148}. La calidad del producto documental se interpreta con el modelo de
ISO/IEC 25010:2023, del que la \S\ref{sec:metricas} deriva las métricas operativas usadas en la
auditoría \cite{iso25010}. El encuadre de procesos de ciclo de vida proviene de ISO/IEC/IEEE
15288:2023 para sistemas y de ISO/IEC/IEEE 12207:2017 para software \cite{iso15288, iso12207}, que
sitúan la definición de requisitos de las partes interesadas y el análisis de requisitos del sistema
dentro de los procesos técnicos.
 
Los fundamentos conceptuales siguen a Pohl \cite{pohl2025} y al SWEBOK v4.0 \cite{swebok4}, y la
lista de verificación de calidad de la especificación se apoya en el temario CPRE Foundation Level
del IREB \cite{ireb31}. La gestión del desarrollo se organiza con Scrum \cite{cohn2004}; el PMBOK
7.ª edición \cite{pmbok7} se emplea de forma complementaria para la planificación y el análisis de
riesgo del proyecto integrador, sin sustituir el marco ágil que rige las iteraciones del equipo.
 
Para el componente de inteligencia artificial de SGCV-IA se consideran dos marcos normativos con
estatus distinto. La Ley Orgánica de Protección de Datos Personales del Ecuador rige de forma
obligatoria el tratamiento de los datos personales de propietarios y los datos clínicos de las
mascotas que el sistema recoge y procesa \cite{lopdp2021}. El Reglamento (UE) 2024/1689 no tiene
vigencia en Ecuador y no se declara como normativa aplicable; se retiene únicamente como marco de
referencia comparativa de buenas prácticas para la clasificación de riesgo y la transparencia de
sistemas de IA \cite{aiact2024}, y así se indica explícitamente cada vez que se cita. Adicionalmente,
por tratarse de un sistema para clínicas veterinarias, se consideró la normativa ecuatoriana
sectorial de ejercicio profesional veterinario y bienestar animal como restricción de dominio sobre
los procesos que el sistema automatiza, aunque no impone requisitos de protección de datos
adicionales a los ya cubiertos por la LOPDP.
 
\subsection{Fundamento teórico}
 
\subsubsection*{Validación frente a verificación}
 
La distinción es anterior a cualquier técnica y condiciona qué se puede concluir de una inspección.
Verificar es comprobar que el documento está bien construido conforme a un estándar; validar es
comprobar que especifica lo que las partes interesadas necesitan \cite{pohl2025}. Un ERS puede
superar la verificación y fallar la validación: sus requisitos serían completos, consistentes y
verificables, y describirían un sistema que la Clínica Veterinaria Macay no pidió.
 
ISO/IEC/IEEE 29148:2018 sitúa la validación en su cláusula 6.4 y enumera las características
exigibles a un requisito individual ---necesario, apropiado, no ambiguo, completo, singular,
factible, verificable, correcto y conforme--- y al conjunto ---completo, consistente, factible,
comprensible, capaz de ser validado \cite{iso29148}. Las métricas de la \S\ref{sec:metricas}
operacionalizan una parte de esas características; las restantes no se midieron porque no admiten un
conteo objetivo sin introducir juicio del auditor.
 
\subsubsection*{Técnicas de validación aplicadas}
 
Las técnicas de validación se ordenan por grado de formalidad: revisión informal, recorrido
(\emph{walkthrough}), inspección formal, prototipado y listas de verificación \cite{pohl2025,
swebok4}. El equipo de SGCV-IA validó el ERS mediante \emph{walkthrough} estructurado, ejecutado al
cierre de cada iteración de Scrum junto con los criterios de aceptación de las historias de usuario
correspondientes, en lugar de una inspección formal tipo Fagan. La razón es de encaje con el proceso:
un ciclo de inspección formal completo ---con preparación individual y reunión de roles
diferenciados--- introduce una cadencia incompatible con \emph{sprints} cortos, mientras que el
\emph{walkthrough} al cierre de iteración permite corregir el catálogo de requisitos antes de que el
error se propague a la siguiente entrega, con un costo de coordinación menor.
 
\subsubsection*{Gestión de requisitos: línea base, trazabilidad y cambio}
 
La gestión de la configuración del ERS descansa en la noción de línea base: una versión aprobada y
congelada que sirve de referencia y a partir de la cual todo cambio se tramita \cite{swebok4}. Sin
línea base no hay control de cambios posible, porque no existe el estado respecto del cual algo
cambia. En SGCV-IA la línea base se materializa como etiquetas de versión sobre el repositorio
Git del proyecto (\texttt{github.com/ramaguas-ship-it/SGCV-IA}), coherentes con las versiones
sucesivas del documento (2.0, con MoSCoW aplicado; 3.0, con los requisitos no funcionales
incorporados).
 
La trazabilidad se distingue en \emph{pre-traceability} ---del requisito hacia su origen--- y
\emph{post-traceability} ---del requisito hacia los artefactos que lo realizan--- \cite{gotel1994}.
Ambas direcciones responden preguntas distintas: la primera, por qué existe un requisito; la segunda,
qué hay que tocar si cambia. La matriz de la \S\ref{sec:gestion} las implementa como columnas hacia
atrás ---hacia la evidencia de elicitación (\id{ENTR-01} a \id{ENTR-08}, encuesta)--- y hacia
adelante ---hacia el modelado, el código y las pruebas--- de una misma cadena.
 
El proceso de cambio sigue el modelo de control integrado de cambios del PMBOK como referencia de
planificación \cite{pmbok7}, adaptado al ritmo de Scrum: la solicitud de cambio se analiza y prioriza
en el refinamiento de \emph{backlog} antes de incorporarse a un \emph{sprint}. La literatura de
trazabilidad advierte de que el análisis de impacto estimado sin recorrer la matriz es la causa más
frecuente de subestimación del costo de un cambio \cite{clelandhuang2012}.
 
\subsubsection*{Herramientas CASE y trazabilidad en entornos ágiles}
 
Las herramientas CASE de gestión de requisitos se clasifican por su posición en el ciclo de vida
---\emph{upper}, \emph{lower}, integradas--- y por el conjunto de funcionalidades esperables:
repositorio, atributos tipados, trazabilidad multinivel, \emph{baselining}, control de cambios,
colaboración e integración con el control de versiones \cite{wiegers2013}.
 
En SGCV-IA la trazabilidad ágil se reconstruye por convención sobre el mensaje de \emph{commit} y la
vinculación de ramas al identificador de requisito o historia \cite{clelandhuang2012}. El formato de
historia de usuario \cite{cohn2004} cumple, dentro del flujo de trabajo diario del equipo, el papel
que en la especificación documental cumple la ficha de requisito conforme a 29148. El proyecto
mantiene ambas representaciones sincronizadas: el ERS como especificación formal de referencia y el
\emph{backlog} de Scrum como su traducción operativa para el desarrollo.
 
\subsubsection*{Ingeniería de requisitos para el componente de inteligencia artificial}
 
Especificar un componente cuyo comportamiento se induce de datos plantea problemas que la ingeniería
de requisitos clásica no aborda. El resultado es probabilístico, de modo que un criterio de
aceptación binario no aplica y hay que sustituirlo por una métrica con umbral y conjunto de
evaluación declarado. La calidad depende del conjunto de datos tanto como del modelo, lo que convierte
las propiedades de los datos ---historiales clínicos, registros de consulta--- en requisitos. Y
aparecen categorías de calidad nuevas ---explicabilidad frente al personal veterinario y ausencia de
sesgo en las recomendaciones--- que ningún modelo de calidad tradicional contemplaba
\cite{vogelsang2024, chazette2021}.
 
A esos requisitos técnicos se suman los normativos: la LOPDP rige el tratamiento de los datos
clínicos y personales que alimentan el modelo de IA de SGCV-IA y los derechos del titular sobre ellos
\cite{lopdp2021}, mientras que el enfoque de clasificación por riesgo del Reglamento (UE) 2024/1689
se usa solo como referencia de buena práctica para documentar el nivel de riesgo del componente de
IA, sin que ello implique una obligación legal en el contexto ecuatoriano \cite{aiact2024}.
 
\subsection{Proceso seguido a lo largo del proyecto}
 
\begin{table}[H]
\centering\footnotesize
\caption{Proceso de ingeniería de requisitos seguido en SGCV-IA}
\label{tab:proceso}
\begin{tabular}{|C{2.6cm}|L{4.0cm}|L{4.6cm}|L{4.6cm}|}
\hline
\rowcolor{grisclaro}\textbf{Fase} & \textbf{Actividad} & \textbf{Técnicas empleadas} & \textbf{Artefacto producido}\\ \hline
Elicitación & Levantamiento de necesidades con la Clínica Veterinaria Macay & Entrevista semiestructurada, encuesta al personal operativo & Entrevistas \id{ENTR-01} a \id{ENTR-08} (ampliadas desde \id{ENTR-09} en la 2B) y encuesta $n=30$, en \texttt{02\_Evidencias/}\\ \hline
Especificación y modelado & Formalización del catálogo de requisitos & Fichas de requisito conforme a 29148, priorización MoSCoW & ERS v2.0 (RF/RNF/RST con MoSCoW) y v3.0 (RNF incorporados)\\ \hline
Validación & Revisión del ERS al cierre de cada iteración & \emph{Walkthrough} estructurado integrado a la revisión de \emph{sprint} & Catálogo de requisitos corregido por iteración\\ \hline
Gestión y trazabilidad & Control de versiones y seguimiento de cambios & Línea base sobre etiquetas de Git, \emph{backlog} de Scrum & Historial de versiones del ERS en el repositorio y matriz de trazabilidad (\S\ref{sec:gestion})\\ \hline
Componente de IA & Especificación de requisitos del módulo de inteligencia artificial & Aplicación de LOPDP y del enfoque de riesgo del Reglamento (UE) 2024/1689 como referencia & \S\ref{sec:ia} del ERS\\ \hline
\end{tabular}
\end{table}
 
\subsection{Proceso de elicitación seguido}
 
La elicitación combinó entrevista semiestructurada con el personal de la Clínica Veterinaria Macay
---ocho entrevistas, codificadas \id{ENTR-01} a \id{ENTR-08} en la versión actual del ERS--- y una
encuesta con $n=30$ aplicada al personal operativo. Para la Entrega 2B el conjunto de entrevistas se
amplía por encima de las ocho ya codificadas, con nuevos identificadores consecutivos a partir de
\id{ENTR-09}, con el fin de cubrir perfiles de la clínica que la primera ronda no alcanzó. La razón de
usar ambos instrumentos y no solo entrevistas es la misma que
justifica la técnica en la literatura: los roles administrativos y clínicos describen el proceso
desde ángulos distintos, y la encuesta permite contrastar a mayor escala lo que las entrevistas
levantan en profundidad. Toda la evidencia de elicitación se conserva anonimizada, con identificador
de analista (\emph{Equipo AHMRV}) y organizada en la carpeta \texttt{02\_Evidencias/} del
repositorio, siguiendo el mismo esquema de codificación con el que se referencia en la matriz de
trazabilidad de la \S\ref{sec:gestion}.
 
Sobre esa base se derivaron los requisitos funcionales (RF), no funcionales (RNF) y restricciones
(RST) del sistema, priorizados con MoSCoW. La versión 2.0 del ERS formalizó el catálogo de requisitos
y aplicó la priorización MoSCoW sobre él; la versión 3.0 incorporó los requisitos no funcionales,
completando el conjunto que la \S\ref{sec:gestion} traza de extremo a extremo.
 
\subsection{Decisiones metodológicas y su justificación}
 
Tres decisiones condicionaron el resultado y merecen justificarse, porque en cada una había una
alternativa razonable que se descartó.
 
\subsubsection*{Elicitación con instrumentos mixtos y no solo con entrevistas}
 
Basar la elicitación solo en entrevistas habría dejado sin contraste cuantitativo lo que el personal
administrativo de la clínica describe como proceso formal frente a lo que el personal operativo
efectivamente hace en la atención diaria. El cuestionario con $n=30$ cumple ese papel de contraste;
las ocho entrevistas semiestructuradas aportan la profundidad que una encuesta cerrada no permite.
Ambos instrumentos se conservan como evidencia trazable e independiente en \texttt{02\_Evidencias/}.
 
\subsubsection*{Priorización con MoSCoW}
 
El catálogo de RF, RNF y RST se priorizó con MoSCoW sobre el conjunto formalizado en la versión 2.0
del ERS. La priorización se mantuvo como criterio único de ordenamiento por alcance del proyecto
integrador: el catálogo resultante es lo bastante acotado para que las categorías \emph{Must},
\emph{Should}, \emph{Could} y \emph{Won't} ordenen por sí solas la secuencia de implementación sin
necesidad de una técnica de puntuación adicional.
 
\subsubsection*{Walkthrough estructurado y no inspección formal}
 
El equipo validó el ERS mediante \emph{walkthrough} al cierre de cada iteración de Scrum, en lugar de
una inspección formal tipo Fagan con roles diferenciados y preparación individual obligatoria
\cite{fagan1976}. La alternativa formal es más rigurosa pero exige una cadencia de reunión y una
preparación previa incompatibles con \emph{sprints} cortos y un equipo reducido; el
\emph{walkthrough} integrado a la revisión de \emph{sprint} permite corregir el catálogo de
requisitos con una frecuencia mayor, al costo de una detección de defectos menos sistemática que la
que ofrecería una inspección formal.
 
\subsubsection*{LOPDP como normativa obligatoria y el Reglamento europeo de IA como referencia}
 
SGCV-IA trata datos personales de los propietarios y datos clínicos de las mascotas, lo que activa de
forma directa la LOPDP ecuatoriana \cite{lopdp2021}. El Reglamento (UE) 2024/1689 no rige en Ecuador;
se decidió, no obstante, mantenerlo como marco de referencia comparativa para documentar el nivel de
riesgo del componente de IA y sus obligaciones de transparencia frente al personal clínico, dejando
explícito en cada referencia que su cumplimiento no es exigible en este contexto sino que se adopta
como buena práctica \cite{aiact2024}.
